\documentclass{article}
\usepackage{amsmath, amssymb, amsthm}
\usepackage{hyperref}
\usepackage{geometry}
\geometry{margin=1in}

\title{Erd\H{o}s Problem \#727}
\date{Last updated: 2025-08-31}
\author{Source: erdosproblems.com}

\begin{document}
\maketitle

\noindent\textbf{Status:} OPEN \quad \textbf{No prize}

\noindent\textbf{Tags:} number theory, factorials

\medskip

\noindent\textbf{Problem Statement:}

Let $k\geq 2$. Does\[(n+k)!^2 \mid (2n)!\]for infinitely many $n$?




A conjecture of Erd\H{o}s, Graham, Ruzsa, and Straus \cite{EGRS75}. It is open even for $k=2$. Balakran \cite{Ba29} proved this holds for $k=1$ - that is, $(n+1)^2\mid \binom{2n}{n}$ for infinitely many $n$. It is a classical fact that $(n+1)\mid \binom{2n}{n}$ for all $n$ (see Catalan numbers ). Erd\H{o}s, Graham, Ruzsa, and Straus observe that the method of Balakran can be further used to prove that there are infinitely many $n$ such that\[(n+k)!(n+1)! \mid (2n)!\](in fact this holds whenever $k&lt;c \log n$ for some small constant $c&gt;0$). Erd\H{o}s \cite{Er68c} proved that if $a!b!\mid n!$ then $a+b\leq n+O(\log n)$.




References

[Ba29] H. Balakran, On the values of $n$ which make $(2n)!/(n+1)!(n+1)!$ an integer . J. Indian Math. Soc. (1929), 97-100.

[EGRS75] Erd\H{o}s, P. and Graham, R. L. and Ruzsa, I. Z. and Straus, E. G., On the prime factors of $(\sp{2n}\sb{n})$ . Math. Comp. (1975), 83-92.

[Er68c] P. Erd\H{o}s, Aufgabe 557 . Elemente Math. (1968), 111-113.

\noindent\textbf{Related OEIS sequences:} A002503, A343507, A389396


\bigskip
\noindent\small{Source: \url{https://www.erdosproblems.com/727}}
\end{document}
